\chapter{Visualization for health data}
\label{ch:visualization}

\begin{quote}
\emph{``The greatest value of a picture is when it forces us to notice what we never expected to see.'' --- John Tukey}
\end{quote}

% ============================================================
\section{Principles of health data visualization}

A good health figure answers a question. A bad one decorates a slide. Before writing any plotting code, ask yourself:
\begin{enumerate}
  \item \textbf{What is the message?} ``Malaria incidence is declining in Sub-Saharan Africa'' --- not ``here is some data.''
  \item \textbf{Who is the audience?} Clinicians need different plots than policymakers.
  \item \textbf{What comparison matters?} Between groups? Over time? Against a threshold?
\end{enumerate}

Rules that apply to every health figure:
\begin{itemize}
  \item Label axes with units (``Systolic BP (mmHg)'', not ``BP'').
  \item Use colorblind-friendly palettes.
  \item Do not use 3D charts, pie charts, or dual y-axes unless you have a very good reason.
  \item Include sample sizes in titles or annotations.
  \item Start y-axes at zero for bar charts. For line plots, zero is optional if the range is narrow.
\end{itemize}

\begin{clinicalnote}
In clinical publications, figures are often the first (and sometimes the only) thing reviewers examine closely. A clear, well-labeled figure can carry a paper. A misleading one can sink it. The Lancet and BMJ both have specific figure guidelines --- study them before submitting.
\end{clinicalnote}

% ============================================================
\section{Setup: matplotlib and seaborn}

\begin{minted}{python}
import pandas as pd
import numpy as np
import matplotlib.pyplot as plt
import seaborn as sns

# Set a clean style for all plots
sns.set_theme(style="whitegrid", font_scale=1.1)
plt.rcParams["figure.figsize"] = (8, 5)
plt.rcParams["figure.dpi"] = 100

# Colorblind-friendly palette
cb_palette = ["#0072B2", "#D55E00", "#009E73",
              "#CC79A7", "#F0E442", "#56B4E9"]
sns.set_palette(cb_palette)
\end{minted}

\begin{pythontip}
The six-color palette above is from Bang Wong's ``Points of view: Color blindness'' (Nature Methods, 2011). It is distinguishable by people with the most common forms of color vision deficiency. Use it as your default.
\end{pythontip}

% ============================================================
\section{Histograms and density plots}

Histograms show the \emph{distribution} of a single variable. In health, this tells you whether a measurement follows a normal distribution, is skewed, or is bimodal.

\begin{minted}{python}
# Load Gapminder data
url = ("https://raw.githubusercontent.com/datasets/"
       "gapminder/main/data/gapminder.csv")
gapminder = pd.read_csv(url)
recent = gapminder[gapminder["year"] == 2007]

# Histogram of life expectancy
fig, axes = plt.subplots(1, 2, figsize=(12, 5))

# Left: histogram
axes[0].hist(recent["lifeExp"], bins=20, edgecolor="white", alpha=0.8)
axes[0].set_xlabel("Life expectancy (years)")
axes[0].set_ylabel("Number of countries")
axes[0].set_title("Distribution of life expectancy, 2007")
axes[0].axvline(recent["lifeExp"].median(), color="red",
                linestyle="--", label=f'Median: {recent["lifeExp"].median():.1f}')
axes[0].legend()

# Right: kernel density estimate (smooth version)
sns.kdeplot(data=recent, x="lifeExp", hue="continent",
            fill=True, alpha=0.3, ax=axes[1])
axes[1].set_xlabel("Life expectancy (years)")
axes[1].set_title("Life expectancy by continent, 2007")

plt.tight_layout()
plt.show()
\end{minted}

\begin{minted}{python}
# Clinical example: distribution of fasting glucose
np.random.seed(42)
normal_glucose = np.random.normal(95, 10, 800)
prediabetic = np.random.normal(115, 8, 150)
diabetic = np.random.normal(180, 40, 50)
all_glucose = np.concatenate([normal_glucose, prediabetic, diabetic])

fig, ax = plt.subplots(figsize=(10, 5))
ax.hist(all_glucose, bins=40, edgecolor="white", alpha=0.7)
ax.axvline(100, color="orange", linestyle="--", linewidth=2,
           label="Normal threshold (100 mg/dL)")
ax.axvline(126, color="red", linestyle="--", linewidth=2,
           label="Diabetic threshold (126 mg/dL)")
ax.set_xlabel("Fasting glucose (mg/dL)")
ax.set_ylabel("Number of patients")
ax.set_title("Distribution of fasting glucose (n=1,000)")
ax.legend()
plt.tight_layout()
plt.show()
\end{minted}

% ============================================================
\section{Box plots}

Box plots compare distributions across groups. The box shows the IQR (25th--75th percentile), the line is the median, and the whiskers extend to 1.5$\times$IQR. Points beyond the whiskers are outliers.

\begin{minted}{python}
# Blood pressure by continent
fig, ax = plt.subplots(figsize=(10, 6))
sns.boxplot(data=recent, x="continent", y="lifeExp", ax=ax)
ax.set_xlabel("Continent")
ax.set_ylabel("Life expectancy (years)")
ax.set_title("Life expectancy by continent, 2007 (n={})".format(len(recent)))
plt.tight_layout()
plt.show()
\end{minted}

\begin{minted}{python}
# Clinical example: BMI by sex and age group
np.random.seed(42)
n = 500
clinical = pd.DataFrame({
    "sex": np.random.choice(["Male", "Female"], n),
    "age_group": np.random.choice(["18-39", "40-59", "60+"], n,
                                  p=[0.3, 0.4, 0.3]),
    "bmi": np.random.normal(27, 5, n).clip(15, 50)
})

fig, ax = plt.subplots(figsize=(10, 6))
sns.boxplot(data=clinical, x="age_group", y="bmi", hue="sex",
            order=["18-39", "40-59", "60+"], ax=ax)
ax.axhline(25, color="orange", linestyle="--", alpha=0.7,
           label="Overweight threshold")
ax.axhline(30, color="red", linestyle="--", alpha=0.7,
           label="Obese threshold")
ax.set_xlabel("Age group")
ax.set_ylabel("BMI (kg/m²)")
ax.set_title(f"BMI by sex and age group (n={n})")
ax.legend()
plt.tight_layout()
plt.show()
\end{minted}

\begin{warningbox}
Box plots hide the shape of the distribution. A bimodal distribution looks the same as a unimodal one in a box plot. Consider using a \textbf{violin plot} (\texttt{sns.violinplot}) or a \textbf{strip plot} (\texttt{sns.stripplot}) overlaid on the box plot for small samples.
\end{warningbox}

% ============================================================
\section{Bar charts}

Bar charts compare \emph{counts} or \emph{rates} across categories. They are the standard choice for disease prevalence comparisons.

\begin{minted}{python}
# Life expectancy comparison: bottom 15 countries (2007)
bottom15 = recent.nsmallest(15, "lifeExp")

fig, ax = plt.subplots(figsize=(10, 7))
ax.barh(bottom15["country"], bottom15["lifeExp"], color="#D55E00")
ax.set_xlabel("Life expectancy (years)")
ax.set_title("15 countries with lowest life expectancy, 2007")
ax.invert_yaxis()  # highest at top
for i, v in enumerate(bottom15["lifeExp"]):
    ax.text(v + 0.5, i, f"{v:.1f}", va="center", fontsize=9)
plt.tight_layout()
plt.show()
\end{minted}

\begin{minted}{python}
# Disease prevalence comparison by country (simulated)
diseases = pd.DataFrame({
    "country": ["Nigeria", "DRC", "Tanzania", "Kenya", "Ghana",
                "South Africa", "Ethiopia", "Uganda"],
    "malaria_prev": [27.3, 31.2, 7.5, 5.8, 15.2, 0.5, 1.2, 9.7],
    "hiv_prev": [1.3, 0.7, 4.7, 4.9, 1.7, 19.0, 1.0, 5.4],
})

fig, ax = plt.subplots(figsize=(10, 6))
x = np.arange(len(diseases))
width = 0.35
ax.bar(x - width/2, diseases["malaria_prev"], width,
       label="Malaria", color="#0072B2")
ax.bar(x + width/2, diseases["hiv_prev"], width,
       label="HIV", color="#D55E00")
ax.set_xticks(x)
ax.set_xticklabels(diseases["country"], rotation=45, ha="right")
ax.set_ylabel("Prevalence (%)")
ax.set_title("Malaria vs HIV prevalence by country")
ax.legend()
plt.tight_layout()
plt.show()
\end{minted}

% ============================================================
\section{Line plots: time trends and epidemic curves}

Line plots show how a measurement changes over time. They are essential for tracking disease trends and epidemic progression.

\subsection{Time trends}

\begin{minted}{python}
# Life expectancy trends by continent
continent_trend = gapminder.groupby(
    ["year", "continent"])["lifeExp"].mean().reset_index()

fig, ax = plt.subplots(figsize=(10, 6))
for continent in continent_trend["continent"].unique():
    subset = continent_trend[continent_trend["continent"] == continent]
    ax.plot(subset["year"], subset["lifeExp"],
            marker="o", markersize=4, label=continent)

ax.set_xlabel("Year")
ax.set_ylabel("Mean life expectancy (years)")
ax.set_title("Life expectancy trends by continent, 1952-2007")
ax.legend(title="Continent")
plt.tight_layout()
plt.show()
\end{minted}

\subsection{Epidemic curves}

\begin{minted}{python}
# Simulate an epidemic curve (SIR-like shape)
np.random.seed(42)
days = pd.date_range("2023-01-01", periods=120, freq="D")
# Simulate daily new cases with an SIR-like peak
t = np.arange(120)
peak = 45
cases = np.maximum(0, 200 * np.exp(-0.5 * ((t - peak) / 15) ** 2)
                   + np.random.normal(0, 10, 120)).astype(int)

epi = pd.DataFrame({"date": days, "new_cases": cases})
epi["rolling_7d"] = epi["new_cases"].rolling(7).mean()
epi["cumulative"] = epi["new_cases"].cumsum()

fig, axes = plt.subplots(2, 1, figsize=(12, 8), sharex=True)

# Top: daily cases + 7-day average
axes[0].bar(epi["date"], epi["new_cases"], alpha=0.4,
            color="#0072B2", label="Daily cases")
axes[0].plot(epi["date"], epi["rolling_7d"], color="#D55E00",
             linewidth=2, label="7-day average")
axes[0].set_ylabel("New cases per day")
axes[0].set_title("Epidemic curve: daily new cases")
axes[0].legend()

# Bottom: cumulative cases
axes[1].plot(epi["date"], epi["cumulative"], color="#009E73",
             linewidth=2)
axes[1].set_xlabel("Date")
axes[1].set_ylabel("Cumulative cases")
axes[1].set_title("Cumulative case count")

plt.tight_layout()
plt.show()
\end{minted}

\begin{clinicalnote}
During outbreaks, the epidemic curve (``epi curve'') is the single most important visualization. Its shape tells you the transmission pattern: a point-source outbreak has a sharp peak; a propagated outbreak has successive waves. The 7-day rolling average smooths out reporting delays and weekend effects.
\end{clinicalnote}

% ============================================================
\section{Scatter plots}

Scatter plots show the relationship between two continuous variables.

\begin{minted}{python}
# GDP per capita vs life expectancy (the classic Gapminder plot)
fig, ax = plt.subplots(figsize=(10, 7))
for continent in recent["continent"].unique():
    subset = recent[recent["continent"] == continent]
    ax.scatter(subset["gdpPercap"], subset["lifeExp"],
               s=subset["pop"] / 1e6,  # bubble size = population
               alpha=0.6, label=continent)

ax.set_xscale("log")
ax.set_xlabel("GDP per capita (log scale, USD)")
ax.set_ylabel("Life expectancy (years)")
ax.set_title("Wealth vs Health, 2007 (bubble size = population)")
ax.legend(title="Continent")
plt.tight_layout()
plt.show()
\end{minted}

\begin{minted}{python}
# Clinical scatter: BMI vs systolic BP with regression line
np.random.seed(42)
n = 200
bmi = np.random.normal(28, 5, n).clip(18, 45)
sbp = 80 + 1.8 * bmi + np.random.normal(0, 12, n)

fig, ax = plt.subplots(figsize=(8, 6))
ax.scatter(bmi, sbp, alpha=0.5, s=30, color="#0072B2")

# Add regression line
from numpy.polynomial.polynomial import polyfit
b, m = polyfit(bmi, sbp, 1)
x_line = np.linspace(bmi.min(), bmi.max(), 100)
ax.plot(x_line, b + m * x_line, color="#D55E00", linewidth=2,
        label=f"y = {m:.1f}x + {b:.1f}")

# Correlation coefficient
from scipy.stats import pearsonr
r, p = pearsonr(bmi, sbp)
ax.set_xlabel("BMI (kg/m²)")
ax.set_ylabel("Systolic BP (mmHg)")
ax.set_title(f"BMI vs Systolic BP (r={r:.2f}, p<0.001, n={n})")
ax.legend()
plt.tight_layout()
plt.show()
\end{minted}

\begin{pythontip}
Seaborn's \texttt{sns.regplot()} and \texttt{sns.lmplot()} add regression lines and confidence bands automatically. Use \texttt{sns.lmplot()} when you want to split by a categorical variable (e.g., regression by sex).
\end{pythontip}

% ============================================================
\section{Kaplan-Meier survival curves}

Survival analysis tracks time until an event (death, relapse, infection). The Kaplan-Meier curve estimates the probability of surviving beyond each time point.

\begin{minted}{python}
# Install lifelines if needed: !pip install lifelines
from lifelines import KaplanMeierFitter

# Simulated survival data: two treatment groups
np.random.seed(42)
n_per_group = 100

# Treatment group: longer survival
treatment_time = np.random.exponential(24, n_per_group).clip(0, 60)
treatment_event = np.random.binomial(1, 0.7, n_per_group)

# Control group: shorter survival
control_time = np.random.exponential(15, n_per_group).clip(0, 60)
control_event = np.random.binomial(1, 0.8, n_per_group)

fig, ax = plt.subplots(figsize=(10, 6))

# Fit and plot treatment group
kmf_treatment = KaplanMeierFitter()
kmf_treatment.fit(treatment_time, event_observed=treatment_event,
                  label="Treatment (n=100)")
kmf_treatment.plot_survival_function(ax=ax, ci_show=True)

# Fit and plot control group
kmf_control = KaplanMeierFitter()
kmf_control.fit(control_time, event_observed=control_event,
                label="Control (n=100)")
kmf_control.plot_survival_function(ax=ax, ci_show=True)

ax.set_xlabel("Time (months)")
ax.set_ylabel("Survival probability")
ax.set_title("Kaplan-Meier survival curves by treatment group")
ax.set_ylim(0, 1.05)

# Add median survival lines
for kmf, color in [(kmf_treatment, "#0072B2"), (kmf_control, "#D55E00")]:
    median = kmf.median_survival_time_
    if not np.isinf(median):
        ax.axhline(0.5, color="gray", linestyle=":", alpha=0.5)
        ax.axvline(median, color=color, linestyle=":", alpha=0.5)

plt.tight_layout()
plt.show()

# Print median survival
print(f"Median survival (treatment): "
      f"{kmf_treatment.median_survival_time_:.1f} months")
print(f"Median survival (control): "
      f"{kmf_control.median_survival_time_:.1f} months")
\end{minted}

\begin{warningbox}
If \texttt{lifelines} is not installed, run \texttt{!pip install lifelines} in your notebook. It is not pre-installed in Google Colab.
\end{warningbox}

\begin{clinicalnote}
Kaplan-Meier curves handle \textbf{censored} observations --- patients who are still alive at the end of the study or who were lost to follow-up. Without proper censoring, you would underestimate survival. The vertical tick marks on the curve show censored observations.
\end{clinicalnote}

% ============================================================
\section{Heatmaps}

Heatmaps display matrices of values using color intensity. Two common uses in health:

\subsection{Correlation matrices}

\begin{minted}{python}
# Correlation between health indicators
np.random.seed(42)
n = 300
health = pd.DataFrame({
    "BMI": np.random.normal(27, 5, n),
    "Systolic_BP": np.random.normal(130, 18, n),
    "Diastolic_BP": np.random.normal(82, 10, n),
    "Glucose": np.random.normal(105, 25, n),
    "Cholesterol": np.random.normal(200, 40, n),
    "HbA1c": np.random.normal(5.8, 1.2, n),
    "Heart_Rate": np.random.normal(72, 12, n),
    "Age": np.random.normal(55, 15, n).clip(18, 90)
})

# Add realistic correlations
health["Systolic_BP"] += 0.8 * health["BMI"]
health["Glucose"] += 15 * health["HbA1c"]
health["Diastolic_BP"] += 0.4 * health["Systolic_BP"]
health["Cholesterol"] += 0.5 * health["BMI"]

corr = health.corr()

fig, ax = plt.subplots(figsize=(9, 8))
mask = np.triu(np.ones_like(corr, dtype=bool))  # hide upper triangle
sns.heatmap(corr, mask=mask, annot=True, fmt=".2f",
            cmap="RdBu_r", center=0, vmin=-1, vmax=1,
            square=True, ax=ax)
ax.set_title("Correlation matrix of health indicators")
plt.tight_layout()
plt.show()
\end{minted}

\subsection{Disease co-occurrence}

\begin{minted}{python}
# Co-occurrence of chronic conditions
conditions = ["Hypertension", "Diabetes", "COPD",
              "Heart Failure", "CKD", "Obesity"]
np.random.seed(42)
cooccurrence = np.random.randint(5, 60, (6, 6))
cooccurrence = (cooccurrence + cooccurrence.T) // 2  # make symmetric
np.fill_diagonal(cooccurrence, [320, 180, 95, 75, 110, 210])

co_df = pd.DataFrame(cooccurrence,
                     index=conditions, columns=conditions)

fig, ax = plt.subplots(figsize=(9, 8))
sns.heatmap(co_df, annot=True, fmt="d", cmap="YlOrRd",
            square=True, ax=ax)
ax.set_title("Disease co-occurrence matrix\n"
             "(diagonal = total cases, off-diagonal = co-occurring)")
plt.tight_layout()
plt.show()
\end{minted}

% ============================================================
\section{Practical: build a 4-panel health dashboard}

Combining multiple plots into a single figure is essential for reports and publications. Here is a complete 4-panel dashboard using real data:

\begin{minted}{python}
import pandas as pd
import numpy as np
import matplotlib.pyplot as plt
import seaborn as sns

sns.set_theme(style="whitegrid", font_scale=1.0)

# Load data
url = ("https://raw.githubusercontent.com/datasets/"
       "gapminder/main/data/gapminder.csv")
df = pd.read_csv(url)
recent = df[df["year"] == 2007].copy()

# Create the 4-panel figure
fig, axes = plt.subplots(2, 2, figsize=(14, 10))
fig.suptitle("Global Health Dashboard, 2007",
             fontsize=16, fontweight="bold", y=1.02)

# ----------------------------------------------------------
# Panel A: Life expectancy distribution by continent
# ----------------------------------------------------------
ax = axes[0, 0]
sns.boxplot(data=recent, x="continent", y="lifeExp", ax=ax,
            palette="Set2")
ax.set_xlabel("")
ax.set_ylabel("Life expectancy (years)")
ax.set_title("A. Life expectancy by continent")
ax.tick_params(axis="x", rotation=30)

# ----------------------------------------------------------
# Panel B: GDP vs Life expectancy scatter
# ----------------------------------------------------------
ax = axes[0, 1]
for continent in recent["continent"].unique():
    subset = recent[recent["continent"] == continent]
    ax.scatter(subset["gdpPercap"], subset["lifeExp"],
               s=subset["pop"] / 2e6, alpha=0.6,
               label=continent)
ax.set_xscale("log")
ax.set_xlabel("GDP per capita (USD, log scale)")
ax.set_ylabel("Life expectancy (years)")
ax.set_title("B. Wealth vs health")
ax.legend(fontsize=8, title="Continent", title_fontsize=8)

# ----------------------------------------------------------
# Panel C: Life expectancy trends over time
# ----------------------------------------------------------
ax = axes[1, 0]
trend = df.groupby(["year", "continent"])["lifeExp"].mean().reset_index()
for continent in trend["continent"].unique():
    subset = trend[trend["continent"] == continent]
    ax.plot(subset["year"], subset["lifeExp"],
            marker="o", markersize=3, label=continent)
ax.set_xlabel("Year")
ax.set_ylabel("Mean life expectancy (years)")
ax.set_title("C. Life expectancy trends, 1952-2007")
ax.legend(fontsize=8, title="Continent", title_fontsize=8)

# ----------------------------------------------------------
# Panel D: Bottom 10 countries bar chart
# ----------------------------------------------------------
ax = axes[1, 1]
bottom10 = recent.nsmallest(10, "lifeExp")
ax.barh(bottom10["country"], bottom10["lifeExp"], color="#D55E00")
ax.set_xlabel("Life expectancy (years)")
ax.set_title("D. 10 lowest life expectancy countries")
ax.invert_yaxis()
for i, v in enumerate(bottom10["lifeExp"]):
    ax.text(v + 0.3, i, f"{v:.1f}", va="center", fontsize=8)

plt.tight_layout()
plt.savefig("health_dashboard_2007.png", dpi=150, bbox_inches="tight")
plt.show()
print("Dashboard saved as health_dashboard_2007.png")
\end{minted}

\begin{exercise}
\begin{enumerate}
  \item Modify the dashboard to show 2002 data instead of 2007. Does the story change?
  \item Add a 5th panel showing a histogram of GDP per capita (log-transformed) with a vertical line at the world median. Use \texttt{fig, axes = plt.subplots(2, 3)} or \texttt{plt.subplot2grid()}.
  \item Create a clinical dashboard for a simulated cohort of 500 patients with the following panels:
    \begin{itemize}
      \item Histogram of BMI with overweight/obese thresholds
      \item Box plot of systolic BP by sex
      \item Scatter plot of age vs.\ HbA1c with a diabetic threshold line at 6.5\%
      \item Bar chart of diagnosis frequency
    \end{itemize}
  \item Using Our World in Data (\url{https://github.com/owid/owid-datasets}), download a dataset on malaria or tuberculosis. Build an epidemic dashboard with:
    \begin{itemize}
      \item Line plot of incidence over time for 5 countries
      \item Bar chart of the most recent year's incidence
      \item Scatter plot of incidence vs.\ GDP
      \item Heatmap of correlation between health indicators
    \end{itemize}
  \item Create a Kaplan-Meier plot comparing survival between 3 treatment groups (drug A, drug B, placebo) using simulated data. Add median survival annotations.
  \item Build a ``pre/post intervention'' visualization: show a health indicator (e.g., malaria cases) as a line plot with a vertical line marking the intervention date. Shade the pre-intervention and post-intervention periods in different colors.
\end{enumerate}
\end{exercise}

% ============================================================
\section{Chapter summary}

\begin{itemize}
  \item Always label axes with units, include sample sizes, and use colorblind-friendly palettes.
  \item \textbf{Histograms} show distributions; add clinical threshold lines to give context.
  \item \textbf{Box plots} compare groups; consider violin plots for more detail.
  \item \textbf{Bar charts} compare counts or rates; use horizontal bars for many categories.
  \item \textbf{Line plots} show trends over time; add 7-day rolling averages for epidemic curves.
  \item \textbf{Scatter plots} reveal relationships; add regression lines and report $r$ values.
  \item \textbf{Kaplan-Meier curves} are the standard for survival data; use the \texttt{lifelines} library.
  \item \textbf{Heatmaps} display correlation matrices and co-occurrence patterns.
  \item Multi-panel figures (\texttt{plt.subplots()}) combine related plots into a coherent story.
  \item Save publication-quality figures with \texttt{plt.savefig("file.png", dpi=300, bbox\_inches="tight")}.
\end{itemize}
